\documentclass[11pt,a4paper]{article}
\usepackage[T1]{fontenc}
\usepackage[utf8]{inputenc}
\usepackage{amsmath,amssymb,amsthm}
\usepackage[margin=1in]{geometry}
\usepackage[colorlinks=true,linkcolor=blue,urlcolor=blue]{hyperref}
\theoremstyle{plain}
\newtheorem{theorem}{Theorem}
\newtheorem{lemma}[theorem]{Lemma}
\newtheorem{proposition}[theorem]{Proposition}
\newtheorem{corollary}[theorem]{Corollary}
\theoremstyle{definition}
\newtheorem{remark}[theorem]{Remark}

\title{The empirical recurrence for OEIS A185499, proved by transfer matrix}
\author{Adrian Perez Fontelles\\ \small Independent researcher}
\date{3 September 2026}
\begin{document}
\maketitle

\begin{abstract}
OEIS A185499 counts arrays over $\{0,\dots,2\}$ under a condition comparing a statistic of
every $2\times2$ subblock with those of its neighbours, and carries an empirical recurrence
of order $73$ contributed by R.~H.~Hardin. It is true, and it is decidable rather than
empirical. The count is a walk count on $S=1944$ states and is $C$-finite; whether the
conjectured recurrence annihilates it is then settled by a finite exact computation.
\end{abstract}

\noindent\small 2020 Mathematics Subject Classification. 05A15, 05B45, 15B36.\normalsize

\section{The conjecture}

OEIS A185499 is ``Number of (n+1)X3 0..2 arrays with every 2X2 subblock diagonal sum less antidiagonal sum equal to some horizontal or vertical neighbor 2X2 subblock diagonal sum less antidiagonal sum.''. It has offset $1$ and begins
\[
99, 697, 4451, 30775, 207705, 1418183, 9658679, 65903901,\ \dots
\]
The entry states, as a conjecture contributed by its author and never marked settled:
\begin{quote}
Empirical: a(n)=16*a(n-1)-66*a(n-2)-134*a(n-3)+1443*a(n-4)-2944*a(n-5)+4970*a(n-6)-14851*a(n-7)-11571*a(n-8)+61044*a(n-9)+159738*a(n-10)+413413*a(n-11)-2914428*a(n-12)-2469514*a(n-13)+6557377*a(n-14)+39589072*a(n-15)+22764483*a(n-16)-206879740*a(n-17)-346733427*a(n-18)+278776495*a(n-19)+1495676084*a(n-20)+2695509571*a(n-21)-4686678083*a(n-22)-7299225497*a(n-23)+6056115083*a(n-24)-16127233866*a(n-25)-88281813514*a(n-26)-7261550680*a(n-27)+1088424262105*a(n-28)+899731654396*a(n-29)-3245433499238*a(n-30)-7815484803768*a(n-31)-406652905248*a(n-32)+22811862007908*a(n-33)+22421132651472*a(n-34)-15498436106752*a(n-35)-12123613528560*a(n-36)-10766636743176*a(n-37)-5067935559232*a(n-38)-408175178563744*a(n-39)-32063168173376*a(n-40)+11604755392032*a(n-41)+674372049197920*a(n-42)+1660897630461888*a(n-43)+3566535713014464*a(n-44)-377394253895680*a(n-45)-1288439207533312*a(n-46)-10038152022954496*a(n-47)-16065995298211328*a(n-48)-17310921165163520*a(n-49)-7179507436965888*a(n-50)+9381015453671424*a(n-51)+65340992207912960*a(n-52)+65285348845109248*a(n-53)+92246407919616000*a(n-54)+129518737725161472*a(n-55)-10289528873484288*a(n-56)-278577260634832896*a(n-57)-222025407375605760*a(n-58)-108195623685586944*a(n-59)-215332645430624256*a(n-60)-241001042326585344*a(n-61)+66507208516435968*a(n-62)+264107131416870912*a(n-63)+164236535305076736*a(n-64)+145690315255185408*a(n-65)+161037241144049664*a(n-66)+51643863923687424*a(n-67)-99781853552050176*a(n-68)-113923742849040384*a(n-69)-40634522378698752*a(n-70)+14258715897102336*a(n-71)+17583870987730944*a(n-72)+5072270477230080*a(n-73)
\end{quote}
As of the ``Last modified'' line on the live entry (Oct 03 2025 03:03:05, revision 11) this is still
recorded as empirical, and nothing on the entry records it as proved.

\section{What is being counted}

Write $\sigma$ for its diagonal sum minus its antidiagonal sum, taken on a $2\times2$ subblock. The contiguous $2\times2$
subblocks of an $(n+1)\times3$ array over $\{0,\dots,2\}$ form a grid of $n$ rows
and $2$ columns, and $\sigma$ labels that grid. The entry
requires that every subblock carries the same value of $\sigma$ as at least ONE of its neighbours horizontally or vertically.

\section{The count is a walk count}

This condition is NOT a condition on a pair. Whether a subblock has a
matching neighbour cannot be decided when its own row is placed: the row BELOW it may yet
supply the match. So the state carries, besides the $2$ consecutive array rows, a tally for
each of the $2$ subblocks of the row just completed --- a single bit saying whether it has already been matched --- counting its matches
among its own row and the row above.

A step appends an array row. It fixes the next subblock row, adds the matches between the two
rows to the OLD row's tallies, and requires those totals to be acceptable, the old row being
final at that moment; the new row's tallies start from its own horizontal matches plus the
matches just made. A state is ACCEPTING when its own tally is already acceptable, which is
exactly the right condition for the last subblock row, the one with nothing below it. Only the
states whose tally came from their own row may BEGIN an array, since the first row has nothing
above it.

An array of $n+1$ rows is then a walk of $n-1$ steps and
\[
a(n)\;=\;\iota^{\!\top}M^{\,n-1}\tau
\]
with $\iota$ the indicator of the starting states and $\tau$ of the accepting ones, on
$S=1944$ states in all. Getting the two vectors right is the whole of the argument: with
$\iota$ taken as all-ones the first row would be allowed a match from a row above it that does
not exist, and the count would be too large. The entry's own first terms settle it --- for
the smallest cases
the model reproduces every published term.

Over the alphabet $\{0,\dots,2\}$ the statistic takes values in $[-4,4]$, so the
grid it labels carries at most $9$ different labels. At $n=1$ the array has $2$
rows and the grid is a single row of $2$
subblocks, with only the neighbours inside that row; the entry's first term is
$99$, which the model reproduces along with every later term, and that is the cheapest
check that the grid has been read with the right shape.

In particular $a$ is $C$-finite. The alignment of the walk index with the entry's own $n$ was
checked against its published terms rather than assumed.

\section{The proof}

\begin{lemma}\label{lem:crit}
Let $q(t)=t^{73}-\sum_i c_it^{73-i}$ be the characteristic polynomial of the
conjectured recurrence. The residual $a(n)-\sum_ic_ia(n-i)$ equals
$\iota^{\!\top}M^{\,j}q(M)\tau$ for the corresponding walk index $j$, so the recurrence
holds from some point on if and only if $\iota^{\!\top}M^{j}q(M)\tau=0$ for all large $j$,
and it is enough to examine $j<S$.
\end{lemma}

\begin{proof}
Factor $M^{j}$ out of $M^{j+73}-\sum_ic_iM^{j+73-i}$. For the bound, $M^{S}$
is an integer combination of $I,M,\dots,M^{S-1}$ by Cayley--Hamilton, so the numbers
$u_j=\iota^{\!\top}M^{j}q(M)\tau$ obey the monic recurrence given by the characteristic
polynomial of $M$; $S$ consecutive zeros force every later one, and the last nonzero $u_j$
pins the threshold exactly.
\end{proof}

\begin{theorem}
\[
a(n)\;=\;16\,a(n-1) - 66\,a(n-2) - 134\,a(n-3) + 1443\,a(n-4) - 2944\,a(n-5) + 4970\,a(n-6) - 14851\,a(n-7) - 11571\,a(n-8) + 61044\,a(n-9) + 159738\,a(n-10) + 413413\,a(n-11) - 2914428\,a(n-12) - 2469514\,a(n-13) + 6557377\,a(n-14) + 39589072\,a(n-15) + 22764483\,a(n-16) - 206879740\,a(n-17) - 346733427\,a(n-18) + 278776495\,a(n-19) + 1495676084\,a(n-20) + 2695509571\,a(n-21) - 4686678083\,a(n-22) - 7299225497\,a(n-23) + 6056115083\,a(n-24) - 16127233866\,a(n-25) - 88281813514\,a(n-26) - 7261550680\,a(n-27) + 1088424262105\,a(n-28) + 899731654396\,a(n-29) - 3245433499238\,a(n-30) - 7815484803768\,a(n-31) - 406652905248\,a(n-32) + 22811862007908\,a(n-33) + 22421132651472\,a(n-34) - 15498436106752\,a(n-35) - 12123613528560\,a(n-36) - 10766636743176\,a(n-37) - 5067935559232\,a(n-38) - 408175178563744\,a(n-39) - 32063168173376\,a(n-40) + 11604755392032\,a(n-41) + 674372049197920\,a(n-42) + 1660897630461888\,a(n-43) + 3566535713014464\,a(n-44) - 377394253895680\,a(n-45) - 1288439207533312\,a(n-46) - 10038152022954496\,a(n-47) - 16065995298211328\,a(n-48) - 17310921165163520\,a(n-49) - 7179507436965888\,a(n-50) + 9381015453671424\,a(n-51) + 65340992207912960\,a(n-52) + 65285348845109248\,a(n-53) + 92246407919616000\,a(n-54) + 129518737725161472\,a(n-55) - 10289528873484288\,a(n-56) - 278577260634832896\,a(n-57) - 222025407375605760\,a(n-58) - 108195623685586944\,a(n-59) - 215332645430624256\,a(n-60) - 241001042326585344\,a(n-61) + 66507208516435968\,a(n-62) + 264107131416870912\,a(n-63) + 164236535305076736\,a(n-64) + 145690315255185408\,a(n-65) + 161037241144049664\,a(n-66) + 51643863923687424\,a(n-67) - 99781853552050176\,a(n-68) - 113923742849040384\,a(n-69) - 40634522378698752\,a(n-70) + 14258715897102336\,a(n-71) + 17583870987730944\,a(n-72) + 5072270477230080\,a(n-73)
\]
for every $n>73$.
\end{theorem}

\begin{proof}
The vector $w=q(M)\tau$ was formed by $73$ matrix--vector products in exact integer
arithmetic and $\iota^{\!\top}M^{j}w$ evaluated for $j=0,1,\dots$, again exactly; those
integers vanish from the index corresponding to $n=73+1$ onwards and the run continued
until the working vector was identically zero, which settles every larger $j$ at once.
\end{proof}

\section{Verification}

Three checks, each independent of the algebra above.

First, the model reproduces all $19$ terms the entry publishes, exactly, in integer
arithmetic. This is what ties the model to the entry: the digraph is built by reading the
entry's English, and a misreading gives different counts.

Second, the conjectured recurrence was evaluated directly on those published terms, with no
matrices involved, and holds wherever the proved range and the data overlap.

Third, the annihilation test of Lemma~\ref{lem:crit} was carried out in exact integer
arithmetic on the whole vector rather than on a sampled prefix.

\begin{thebibliography}{9}
\bibitem{oeis} The OEIS Foundation, \emph{The On-Line Encyclopedia of Integer
Sequences}, \url{https://oeis.org}, 2026. Sequence A185499.
\bibitem{stanley} R.~P.~Stanley, \emph{Enumerative Combinatorics}, Volume 1, 2nd ed.,
Cambridge University Press, 2011. (Transfer-matrix method, Section 4.7.)
\bibitem{hu} J.~E.~Hopcroft, R.~Motwani and J.~D.~Ullman, \emph{Introduction to Automata
Theory, Languages, and Computation}, 3rd ed., Addison--Wesley, 2006.
\bibitem{horn} R.~A.~Horn and C.~R.~Johnson, \emph{Matrix Analysis}, 2nd ed., Cambridge
University Press, 2013.
\end{thebibliography}

\end{document}
